\section{DLat (Latices Distributivos)}

Os latices distributivos são estruturas onde qualquer par de elementos possui um ínfimo e um supremo, e onde estas operações distribuem-se entre si. São fundamentais para o estudo de álgebras de lógica e conjuntos.

\subsection{Objetos}
Utilizamos uma abordagem de programação funcional em MATLAB, recorrendo a estruturas que armazenam os "handles" das operações de \textit{meet} ($\wedge$) e \textit{join} ($\vee$).

\begin{lstlisting}[caption={Estrutura Funcional de Latice}]
L.elements = {1, 2, 3, 4}; 
L.meet = @(a, b) min(a, b); 
L.join = @(a, b) max(a, b); 
\end{lstlisting}

\subsection{Morfismos}
Um morfismo em \textbf{DLat} é uma função que preserva ambas as operações binárias. A validação exige testar a igualdade $f(a \wedge b) = f(a) \wedge f(b)$ para todos os pares do domínio.